\documentclass[../main.tex]{subfiles}

\begin{document}

\chapter{集合论与逻辑}\label{chapter:ch01-集合论与逻辑}

我们像大多数数学家那样，对于集合论采取一种朴素的观点。我们认定由一些对象构成的集合这个概念是直观自明的，本书在此基础上展开讨论，而不去进一步分析这个概念。有关集合的深入研究属于数学基础和数理逻辑的范畴，而研究这些领域不是我们的目的。

逻辑学家十分详细地分析过集合论，并且就此论题总结出了若干公理。每一个这样的公理都阐述被数学家所普遍接受的集合的一个性质，这些公理为数学其他领域奠定了广泛和坚实的基础。

不幸的是，仅凭直觉而粗心地使用集合论难免会导致谬误。事实上，将集合论公理化的原因之一，便在于建立一些处理集合的相关法则以避免造成这些谬误。尽管我们并不深入地讨论这些公理，但我们所遵从的那些处理集合的法则都源于这些公理。在本书中，读者将会学到如何按“适当”的方式对待集合：体察我们的处理方式，再加上读者自己的解题实践。或许你出于学习的需要，希望对集合论有更深入、更详细的了解，那么学一门逻辑学课程或数学基础课程也许是一个更好的选择。

\subfile{Chapter-01/section-01.tex}
\subfile{Chapter-01/section-02.tex}
\subfile{Chapter-01/section-03.tex}
\subfile{Chapter-01/section-04.tex}
\subfile{Chapter-01/section-05.tex}
\subfile{Chapter-01/section-06.tex}
\subfile{Chapter-01/section-07.tex}
\subfile{Chapter-01/section-08.tex}
\subfile{Chapter-01/section-09.tex}
\subfile{Chapter-01/section-10.tex}
\subfile{Chapter-01/section-11.tex}
\subfile{Chapter-01/section-extra.tex}

\end{document}